\section{Finding-by-finding placement in Grothendieck's frameworks}
\label{sec:framework-map}

The following map locates the framework of each group of findings. It
does not attribute the new constants, minima, algorithms, or counterexamples
to the historical texts. The exact mathematical input used in a proof is
cited at its point of use; an archive passage here gives its conceptual
location and the relevant distinction of equivalence relations. In
particular none of the numerical bounds is read out of a conjectural
archive assertion.

\small
\begin{longtable}{@{}p{.19\textwidth}p{.35\textwidth}p{.39\textwidth}@{}}
\toprule
Findings & Geometric or motivic statement & Exact framework location\\
\midrule
\endfirsthead
\toprule
Findings & Geometric or motivic statement & Exact framework location\\
\midrule
\endhead
B01--B02, B02a,\newline B05--B06 & Bounded integral cycles, their parameter spaces, and equality in a fixed realization; descent and the coefficient correction. & FGA, Bourbaki Expos\'e221, Lemma2.4, p.255, and Theorem3.2, pp.259--260; SGA$4\frac12$, \emph{Cycle}, 2.3.1--2.3.3, 2.3.7--2.3.8.\\[4pt]
B03--B04 & Specialization of homological classes and a degree-dependent generic-flatness open set. & EGA IV$_2$, Theorem6.9.1, p.154; SGA4, XVI, Corollary2.2, p.214, with the locally constant coefficient hypothesis.\\[4pt]
C01--C07, C18 & Rational correspondences, realization-compatible CM transport, and effective representatives after translated pullback. & Cote16, scans48--49: algebraic inverses and the distinction between an isomorphism and a split monomorphism; cote14, scan9: effectivity cones. The archive provides no quantitative CM-transport theorem.\\[4pt]
P06--P14; V01--V07 & Projectors onto a summand of the homological motive; positive intersection forms; exact integral representative costs and their failure to depend on the norm alone. & Cote14, scans44--45: algebraic versus homological formulations and the primitive sign; cote16, scan49: correspondence images and the index relation. The explicit minima and saturated lattice are new calculations.\\[4pt]
P02--P03; C30--C32 & Frobenius, geometric endomorphisms, and compatibility of a CM marking with reduction. & SGA4, XVI, Corollary2.2 supplies the specialization framework. The elliptic prime classification and quaternion-order calculations use the separate cited elliptic-curve theorems; they are not assertions of this SGA corollary.\\[4pt]
W01--W04 & Effective classes and homological replacement by sums of abelian subvarieties, with integral multiplicities retained; an exact finite class comparison. & Cote14, scan9: cone/effectivity notation; scans44--45: the equivalence relation must be stated. These findings assert homological equality, not algebraic equivalence.\\[4pt]
R10--R12; R20--R24 & Algebraic projector extraction, inverse intersection matrices, charged isogenies, and contraction of correspondences. & Cote16, scans48--49: algebraic inverse correspondences and the extra index hypothesis for splitting; cote14, scan45: signed primitive positivity. FGA Expos\'e221, Theorem3.2 supplies the proper parameter spaces used for extending bounded detection.\\[4pt]
R13a--R13c & A finite cover trivializing the Weil local system, closed degree thresholds, and a family bound. & SGA1, XII, Theorem5.1 for algebraizing the finite cover, and Proposition2.4 for analytic/algebraic connectedness; FGA Expos\'e221, Theorem3.2 for the proper parameter spaces.\\[4pt]
R30--R31 & A finite-prime test for the generic-flatness gate. & EGA IV$_2$, Theorem6.9.1 supplies that gate. The numerical prime range comes from the separately cited Dirichlet-GRH estimate, not from EGA or an archive conjecture.\\[4pt]
R40--R43 & Exact cycle-class comparison in Tate or Betti realizations, using finite bounded parameter spaces. & FGA Expos\'e221, Lemma2.4 and Theorem3.2; SGA$4\frac12$, \emph{Cycle}, 2.3.1--2.3.8. The comparison with numerical intersection vectors additionally uses the stated Tate and equivalence inputs.\\[4pt]
R44--R45 & An explicit generic-flatness certificate and a computable finite exceptional table. & EGA IV$_2$, Theorem6.9.1, p.154, made effective here by tracked monic bases on affine charts.\\[4pt]
C15--C16 & A primitive-sign correction and an audit of what qualitative source theorems actually bound. & Cote14, scan45 includes the sign $(-1)^m$; cote16, scan49 distinguishes positivity-assisted splitting from effective-cycle control. C15 concerns Milne's supplied report, not a transcription error in Grothendieck.\\
\bottomrule
\end{longtable}
\normalsize

All findings in this paper concern algebraic cycles modulo homological
equivalence unless a different relation is explicitly stated. When a
numerical intersection test is used, its equality with the relevant
homological test is an explicit theorem input. No passage from
homological equality to algebraic equivalence is implicit in the
terminology of motives or correspondences.
